\chapter{数式データの扱い}\label{mexpr}

本章は，シンボリックな計算を扱うためのEgisonの機能を紹介する．

\section{数式データの内部表現}\label{mexpr-expr}

数式は組み込みデータ型としてEgison処理系内部でHaskellの代数的データ型を用いて直接実装されている．
\lstinline{fromMathExpr}という組み込み関数を使うと，この数式データをEgison上の代数的データ型に変換することができる．
\begin{lstlisting}[language=egison,numbers=none]
fromMathExpr  (x - sqrt 2)
-- Div (Plus [Term 1 [(Symbol x [], 1)], Term -1 [(Apply sqrt [Div (Plus [Term 2 []])
--     (Plus [Term 1 []])], 1)]]) (Plus [Term 1 []])
\end{lstlisting}
本節は，数式データを表現するこの内部の代数的データ型について説明する．

上記の出力の代数的データ型からもわかるように，数式データ型は，分母と分子の両方に積和標準形の多項式をとる分数として定義されている．
\lstinline{Div}がこの分数を生成するためのデータコンストラクタで，\lstinline{Plus}が多項式を生成するためのデータコンストラクタとなっている．
そのため，たとえば分数同士の足し算の結果は一つの分数にまとめられる．
\begin{lstlisting}[language=egison,numbers=none]
a / b + x / y
-- (a * y + x * b) / (b * y)
\end{lstlisting}
\lstinline{Plus}の引数がリストであることからもわかるように，多項式は，項のリストからなるデータ型として定義されている．
そして，項は係数と因子のリストからなるデータ型として定義されている．
項は\lstinline{Term}データコンストラクタを使って生成される．
\lstinline{Term}は第一引数に係数，第二引数に因子のリストをとる．
係数には整数，因子にはシンボルなどをとることができる．
たとえば，\lstinline{3 * x^2 * sqrt 2}という項は，係数として\lstinline{3}，因子のリストとして\lstinline{[x^2, sqrt 2]}をとる．
因子にはシンボル以外に，$\sqrt{2}$のような項もとることもできる．
因子のリストは，底（$x^2$の場合$x$）と冪指数（$x^2$の場合$2$）の連想リストとなっている．
冪指数には自然数しかとれないという制限がある．
\begin{lstlisting}[language=egison,numbers=none]
fromMathExpr (3 * x^2 * sqrt 2)
-- Div (Plus [Term 3 [(Symbol x [], 2), (Apply sqrt [Div (Plus [Term 2 []]) (Plus [Term 1 []])], 1)]]) (Plus [Term 1 []])
\end{lstlisting}
上記の例のように，因子は\lstinline{Symbol}や\lstinline{Apply}といったデータコンストラクタを使って生成される．
\lstinline{Symbol}データコンストラクタは，第一引数にシンボル自身を，第二引数に添字のリストをとる．
\lstinline{Symbol}データコンストラクタの第一引数にシンボル自身にしているのは，\lstinline{fromMathExpr}の結果を扱いやすくするためである．
シンボルを表現するのデータ構造のなかにシンボル自身があると定義が無限に循環してしまうため，Haskellによるシンボルの内部表現は\lstinline{Symbol}データコンストラクタの第一引数は文字列をとっている．
添字については，テンソルについて解説する第\ref{tensor}章で説明する．
\lstinline{Apply}データコンストラクタは，第一引数に関数名のシンボルを，第二引数に引数の数式のリストをとる．
その他に因子のデータコンストラクタに\lstinline{Quote}と\lstinline{Function}がある．
\lstinline{Quote}データコンストラクタについては，\ref{mexpr-backquote}節で紹介する．
\lstinline{Function}データコンストラクタについては，第\ref{funcsym}章で紹介する．

\lstinline{fromMathExpr}の逆にEgison上の代数的データ型として表現された数式をEgison処理系内部の数式データに変換する\lstinline{toMathExpr}関数も用意されている．
\lstinline{toMathExpr}関数を使うと，Egison上の代数的データ型として数式を処理したあとに，これらのデータをEgison内部の数式データに戻すことができる．
\begin{lstlisting}[language=egison,numbers=none]
toMathExpr (fromMathExpr  (x - sqrt 2))
-- x - sqrt 2
\end{lstlisting}
ただし，\lstinline{fromMathExpr}と\lstinline{toMathExpr}のような組み込み関数は，Egisonライブラリで定義されている\lstinline{mathExpr}マッチャーを使えば，ユーザーが触る必要はほぼない．
\lstinline{mathExpr}マッチャーの実装に\lstinline{fromMathExpr}と\lstinline{toMathExpr}は使われている．

数は複雑なデータ型である．
自然数までであれば代数的データ型としてエレガントに定義できる．
しかし，自然数の足し算の逆関数引き算を考えると，負の数という概念が得られる．
さらに足し算を繰り返すことによって得られる掛け算の逆を考えると，有理数という概念が得られる．
さらに掛け算を繰り返すことによって得られるべき乗という操作の逆を考えると，平方根や虚数，それらを含む代数的数の概念が得られる．

\section{バッククオート（\texttt{`}）による式展開の制御}\label{mexpr-backquote}

Egisonは数式を自動で積和標準形に展開すると\ref{mexpr-symbol}で述べたが，この展開をバッククオート（\lstinline{`}）により制御することができる．
バッククオートに続く式は，ひとつのシンボルのように扱われる．
\begin{lstlisting}[language=egison,numbers=none,escapechar=']
`(x + 1)^2
--  `(x + 1)^2

`(x + 1) + `(x + 1)
--  2 * `(x + 1)

`(x + 1) + (x + 1)
-- `(x + 1) + x + 1
\end{lstlisting}

バッククオートが役に立つのは，主に下記のような多項式の割り算を含む計算のときである．
\begin{lstlisting}[language=egison,numbers=none,escapechar=']
(x + 1)^2 / (x + 1)
-- (x^2 + 2 * x + 1) / (x + 1)

`(x + 1)^2 / `(x + 1)p
-- '(x + 1)
\end{lstlisting}

バッククオートによりクオートされた項は，Egison内部の数式データのなかで因子として扱われる．
\lstinline{Quote}は，\lstinline{Symbol}や\lstinline{Apply}と同列のデータコンストラクタとして定義されている．
\begin{lstlisting}[language=egison,numbers=none,escapechar=']
fromMathExpr (`(x + 1)^2)
-- Div (Plus [Term 1 [(Quote (Div (Plus [Term 1 [(Symbol x [], 1)], Term 1 []]) (Plus [Term 1 []])), 2)]]) (Plus [Term 1 []])fromMathExpr `(x + 1)^2
\end{lstlisting}

\section{シングルクオート（\texttt{'}）による関数適用の制御}\label{mexpr-squote}

Egisonはシンボルを関数として適用するとシンボリックな結果を返す．
例えば，未定義の変数\lstinline{f}についての関数適用は以下のように動作する．
\begin{lstlisting}[language=egison,numbers=none]
f a -- f a
\end{lstlisting}
定義済みの関数の適用についても上記のようにシンボリックな結果を返したいことがある．
たとえば，\lstinline{sqrt}関数については以下のように動作してほしい．
\begin{lstlisting}[language=egison,numbers=none]
sqrt 4 -- 2
sqrt 5 -- sqrt 5
\end{lstlisting}
このような動作を実現するために，Egisonには組み込み構文としてシングルクオート（\lstinline{'}）が用意されている．
シングルクオートには変数が続き，たとえその変数が定義済みであったとしても，その変数の名前のシンボルを返す．
シングルクオートを関数の前に付加すれば，定義済みの関数に対してもシンボリックな関数適用を表現できる．
\begin{lstlisting}[language=egison,numbers=none]
def f x := x + 1

f a -- a + 1
'f a -- f a
\end{lstlisting}
シングルクオートは平方根を計算する\lstinline{sqrt}の実装以外にも，三角関数\lstinline{sin}や\lstinline{cos}の実装などでEgisonライブラリ内で使われている．

\begin{lstlisting}[language=egison,numbers=none]
fromMathExpr x
-- Div (Plus [Term 1 [(Symbol x [], 1)]]) (Plus [Term 1 []])
fromMathExpr 'x
-- Div (Plus [Term 1 [(Symbol x [], 1)]]) (Plus [Term 1 []])
\end{lstlisting}

\section{\texttt{withSymbols}式によるローカルシンボルの宣言}\label{mexpr-withsym}

未定義の変数はシンボルとして扱われるという仕様は，すでに定義済みの変数をユーザーがシンボルとして扱おうとするミスを誘発しやすい．
この問題を解決するために，\lstinline{withSymbols}式というローカルシンボルを宣言するための構文がEgisonには用意されている．
\lstinline{withSymbols}式は，第一引数に変数のリストをとり，これらの変数は第二引数の式の評価時にシンボルとして扱われる．
\begin{lstlisting}[language=egison,numbers=none]
def i := 10

i -- 10
withSymbols [i] i -- i
i + withSymbols [i] i -- 10 + i
\end{lstlisting}

\lstinline{withSymbols}式で宣言されたローカルなシンボルは，未定義変数から生成されるグローバルなシンボルとは独立したシンボルとして処理される．
\begin{lstlisting}[language=egison,numbers=none]
j + withSymbols [j] j -- j + j
\end{lstlisting}

\section{数式データに対するパターンマッチ}\label{mexpr-pattern}

本節は，数式データに対して用意されているパターンを紹介する．
数式に対するマッチャー\lstinline{mathExpr}はEgisonライブラリに実装されている．
\lstinline{mathExpr}の定義は\texttt{lib/math/expression.egi}で確認できる．
また数式に対するパターンを，より数式に近いかたちで記述するためにいくつかの糖衣構文が組み込みで実装されている．
本節では，これらの糖衣構文も紹介する．
数式に対するパターンマッチは，\ref{cas-derivative}節で実演されている．
パターンの具体的な使用例として，\ref{cas-derivative}節のプログラムを参照してほしい．

\subsection{因子に対するパターンマッチ}

因子に対するパターンコンストラクタは，\lstinline{symbol}・\lstinline{apply}・\lstinline{quote}・\lstinline{func}の4つがある．
\lstinline{symbol}はシンボルに，\lstinline{apply}は\ref{cas-symbol2}節で紹介したシンボリックな関数適用に，\lstinline{quote}は\ref{mexpr-backquote}節で紹介したバッククオートされた式に，\lstinline{func}は第\ref{funcsym}章で解説する関数シンボルと呼ばれるオブジェクトにパターンマッチするためのパターンコンストラクタである．

\lstinline{symbol}パターンコンストラクタの第一引数はシンボル自身に，第二引数は添字にパターンマッチする．
添字については，テンソルについて解説する第\ref{tensor}章で改めて紹介する．
\begin{lstlisting}[language=egison,numbers=none]
match x as mathExpr with
  | symbol $s _ -> s
-- x
\end{lstlisting}
シンボルに対しては\lstinline{symbol}パターンコンストラクタを使ってパターンマッチすることもできるが，実際には値パターンを使ってパターンマッチすることが多い．
\begin{lstlisting}[language=egison,numbers=none]
match x as mathExpr with
  | #x -> "Matched"
  | _ -> "Not matched"
-- "Matched"
\end{lstlisting}

\lstinline{symbol}パターンコンストラクタの第一引数は関数名のシンボルに，第二引数は引数の数式のリストにパターンマッチする．
引数のパターンマッチには，\lstinline{list mathExpr}マッチャーが使われるため，リストに対するパターンを使うことができる．
\begin{lstlisting}[language=egison,numbers=none]
match ('f x 1 y)  as mathExpr with
  | apply $f ($x :: $xs) -> (f, x, xs)
-- (f, x, [1, y])
\end{lstlisting}

\todo{quoteパターン}

\subsection{項に対するパターンマッチ}

項に対するパターンマッチには，\lstinline{term}パターンコンストラクタを使う．
\lstinline{term}パターンコンストラクタの第一引数は係数に，第二引数は因子の連想リストにパターンマッチする．
因子の連想リストのパターンマッチには，assocMultisetマッチャーが使われる．\todo{assocMultisetマッチャーの説明}
\begin{lstlisting}[language=egison,numbers=none]
matchAll 3 * x^2 * y  as mathExpr with
  | term $a ($x :: $xs) -> (a, x, xs)
-- [(3, x, [(x, 1), (y, 1)]), (3, y, [(x, 2)])]
\end{lstlisting}

\lstinline{term}パターンコンストラクタには，より数式に近いパターンを記述するための糖衣構文がある．
\begin{lstlisting}[language=egison,numbers=none]
matchAll 3 * x^2 * y  as mathExpr with
  | $a * $x * $xs -> (a, x, xs)
-- [(3, x, x * y), (3, y, x^2)]
\end{lstlisting}

\todo{べき乗パターンの説明}
\begin{lstlisting}[language=egison,numbers=none]
matchAll 3 * x^2 * y  as mathExpr with
  | $a * $x^#2 * $xs -> (a, x, xs)
-- [(3, x, y)]
\end{lstlisting}

\subsection{多項式に対するパターンマッチ}

有理式に対するパターンマッチには，\lstinline{poly}パターンコンストラクタを使う．
\lstinline{poly}パターンコンストラクタは引数に項のリストにマッチするパターンをとる．
\begin{lstlisting}[language=egison,numbers=none]
matchAll x + y + 1  as mathExpr with
  | poly ($x :: $xs) -> (x, xs)
-- [(x, [y, 1]), (y, [x, 1]), (1, [x, y])]
\end{lstlisting}

\lstinline{poly}パターンコンストラクタにも，より数式に近いパターンを記述するための糖衣構文がある．
\begin{lstlisting}[language=egison,numbers=none]
matchAll x + y + 1  as mathExpr with
  | $x + $xs -> (x, xs)
-- [(x, y + 1), (y, x + 1), (1, x + y)]
\end{lstlisting}

\subsection{有理式に対するパターンマッチ}

有理式に対するパターンマッチには，\lstinline{div}パターンコンストラクタを使う．
\lstinline{div}パターンコンストラクタの第一引数は分子の多項式に，第二引数は分母の多項式にパターンマッチする．
\begin{lstlisting}[language=egison,numbers=none]
matchAll (a + b) / c  as mathExpr with
  | div $x $y -> (x, y)
-- [(a + b, c)]
\end{lstlisting}

\lstinline{div}パターンコンストラクタにも，より数式に近いパターンを記述するための糖衣構文がある．
\begin{lstlisting}[language=egison,numbers=none]
matchAll (a + b) / c  as mathExpr with
  | $x / $y -> (x, y)
-- [(a + b, c)]
\end{lstlisting}
